\begin{textsample}{POS Dim 3 – rr – Score 16.00 – rr/rr\_inf\_11.txt}  \label{ex:f3_pos_227}
3.4.4 \textbf{ESCUTA}, FALA, PENSAMENTO E IMAGINAÇÃO Desde o \textbf{nascimento}, as \textbf{crianças} participam de situações comunicativas cotidianas com as pessoas com as quais interagem.
As primeiras formas de interação do \textbf{bebê} são os movimentos do seu corpo, o olhar, a postura corporal, o sorriso, o choro e outros recursos vocais, que ganham sentido com a interpretação do outro.
Progressivamente, as \textbf{crianças} vão ampliando e enriquecendo seu vocabulário e demais recursos de expressão e de compreensão, apropriando -se da língua materna que se torna, pouco a pouco, seu veículo privilegiado de interação.
Por meio da linguagem, nos fazemos presentes, interagimos, aprendemos, \textbf{inventamos}, participamos, argumentamos, narramos.
A linguagem tem também papel fundamental na constituição do pensamento e, consequentemente, na formação das subjetividades, nos constituindo como \textbf{seres} humanos, participantes da cultura na qual estamos inseridos.
Para Vygotsky ( 1996 ), a linguagem é matéria prima para o pensamento.
Segundo ele, em suas origens, pensamento e linguagem possuem raízes diferenciadas e linhas diferentes de evolução.
Num determinado momento do desenvolvimento humano, as linhas do pensamento e da linguagem se cruzam, possibilitando que a linguagem se torne intelectual e o pensamento verbal.
Na Educação \textbf{Infantil}, uma vez que a linguagem é mediadora das diversas interações ocorridas, quer seja entre \textbf{adultos} e \textbf{crianças}, quer entre as próprias \textbf{crianças}, é fundamental propiciar experiências ricas e significativas de fala e \textbf{escuta}, desde os \textbf{bebês}.
Desde cedo, a \textbf{criança} manifesta \textbf{curiosidade} com relação à cultura escrita : ao \textbf{ouvir} e acompanhar a leitura de textos, ao observar os muitos textos que circulam no contexto familiar, comunitário e escolar, ela vai construindo sua concepção de língua escrita, reconhecendo diferentes usos sociais da escrita, dos gêneros, suportes e portadores.
Na Educação \textbf{Infantil}, a \textbf{imersão} na cultura escrita deve partir do que as \textbf{crianças} conhecem e das \textbf{curiosidades} que deixam transparecer.
As experiências com a literatura \textbf{infantil}, propostas pelo educador, mediador entre os textos e as \textbf{crianças}, contribuem para o desenvolvimento do gosto pela leitura, do \textbf{estímulo} à imaginação e da ampliação do conhecimento de mundo.
Além disso, o contato com histórias, contos, fábulas, poemas, cordéis etc. propicia a familiaridade com \textbf{livros}, com diferentes gêneros literários, a diferenciação entre ilustrações e escrita, a aprendizagem da direção da escrita e as formas corretas de manipulação de \textbf{livros}.
Na \textbf{pequena} \textbf{infância} a aquisição e o domínio da linguagem verbal estão vinculados à constituição do pensamento, à fruição literária também é um instrumento de apropriação dos demais conhecimentos.
Nesse convívio com textos escritos, as \textbf{crianças} vão construindo hipóteses sobre a escrita que se revelam, inicialmente, em rabiscos e garatujas e, à medida que vão conhecendo letras, em escritas espontâneas, não convencionais, mas já \textbf{indicativas} da compreensão da escrita como sistema de representação da língua.

% matched lemmas: adulto, bebê, criança, curiosidade, escuta, estímulo, imersão, indicativo, infantil, infância, inventar, livro, nascimento, ouvir, pequeno, seres
\end{textsample}
